\section{The Glory Takes the Road: Overshadowing and the Hill Country}

Israel had known the shock of glory arriving.  At Sinai the cloud descended;
at the completed tabernacle the cloud covered the tent, the glory filled it,
and even Moses could not enter (Exod 40:34--35).  The Ark stood within that
sanctuary as the covenant throne and appointed sign of the Lord's presence.
Centuries later Gabriel comes not to a royal court or a rebuilt Holy of Holies,
but to a virgin in Nazareth.  His promise is at once gentler and more
astonishing: the Holy Spirit will come upon her, the power of the Most High
will overshadow her, and the child born from her will be holy, the Son of God
(Luke 1:26--35).  The glory is no longer approaching a chest within a tent.
The eternal Son is taking flesh from a woman who believes.

The difference makes the fulfillment greater.  Mary is not wood, gold, or a
passive chamber.  She hears, asks, receives an answer, and freely consents that
God's word be fulfilled in her (Luke 1:38).  Nor is the child a
new sacred object.  He is the divine person in whom God and humanity are
united, the Lord whose human nature is truly taken from his mother.  The old
Ark could bear signs of covenant presence.  Mary bears the incarnate Presence
himself.

\subsection{The cloud and the overshadowing}

The Greek Exodus gives the correspondence a precise sound.  After the Ark and
the other furnishings have been placed in the sanctuary, the cloud ``was
overshadowing'' the tent: the verb is \emph{episkiazō}, in the imperfect tense
(Exod 40:35 LXX).  Gabriel uses the same verb in the future tense: the power of
the Most High ``will overshadow'' Mary (Luke 1:35).  The New Testament also
uses the verb for the Transfiguration cloud and for Peter's passing shadow
(Matt 17:5; Mark 9:7; Luke 9:34; Acts 5:15).  It is therefore not a secret word
found only beside the Ark, but its Exodus setting makes the sanctuary resonance
real.

Luke himself confirms the register.  At the Transfiguration the cloud
overshadows the disciples and the Father's voice identifies the beloved Son
(Luke 9:34--35).  Overshadowing in Luke can therefore mark divine presence and
revelation, not biological agency described in pagan terms.  The Holy Spirit
comes upon Mary; the Father's power overshadows; the Son assumes from her the
humanity in which he will save.  The scene is wholly Trinitarian and wholly
graced.

The exact Greek wording ``will overshadow you'' also appears in Psalm 90:4
(91:4 in common English numbering), where God shelters the faithful beneath
his wings.  Luke's sentence can carry both biblical resonances: the
sanctuary-cloud of Exodus and the protecting wings of the psalm.  These senses
belong together.  The power that makes Mary the living sanctuary is not
violence against her person; it is the holy, sheltering action of God to which
her freedom answers in faith.

The Hebrew Exodus keeps the picture concrete: the cloud ``covered'' the tent
and ``settled'' or ``dwelt'' upon it.  It does not use a Hebrew verb meaning
``overshadow,'' and it does not contain the later noun \emph{Shekinah} in this
passage.  The typology does not need either claim.  The inspired scenes already
give what Christian reading requires: cloud, indwelling glory, appointed
sanctuary, and then the Lucan promise that the Most High will overshadow the
woman from whom his Son is born.

\subsection{Mary rises}

The Annunciation does not end in stillness.  Mary rises and goes with haste
into Judah's hill country (Luke 1:39).  Luke records no self-dramatizing speech
on the road.  The mystery travels hidden from human sight: the Lord of Israel
is borne beneath his mother's heart.  The old journeys of the Ark begin to
sound again---Judah's uplands, a holy bearer entering a house, the eruption of
joy, and the canonical theme of blessing at a holy arrival.

At Zechariah's threshold Mary's greeting is heard.  John leaps in Elizabeth's
womb; Elizabeth is filled with the Holy Spirit; and a great cry bursts from
her (Luke 1:40--42).  Nothing in the narrative treats Mary as an independent
source of sanctification.  Christ is present because she bears him, the Spirit
fills Elizabeth, and the child responds with joy.  Yet Mary is not incidental.
Elizabeth's Spirit-filled question names her relation to the hidden Lord: she
wonders how the mother of her Lord should come to her (Luke 1:43).  The grace
belongs to God; the real maternal coming belongs to Mary.

David had once danced before the arriving Ark (2 Sam 6:14--16).  Now the
unborn forerunner's body becomes praise before the mother of his Lord.  The
correspondence is first an image rather than an identity of wording: the
standard Septuagint describes David as dancing and playing, while Luke says
that John leaped in exultation.  The bodily joy nevertheless has one center.
David rejoices before the Lord's Ark; John rejoices before the Lord whom Mary
carries.

\subsection{The question and the three months}

David's wonder had taken the form of a question: how could the Ark of the Lord
come to him?  Elizabeth's wonder takes a closely corresponding form: how
should the mother of her Lord come to her? (2 Sam 6:9; Luke 1:43).  She does
not call Mary the Ark, and the sentences are not a word-for-word quotation.
But within the gathering pattern the question sounds like recognition.  The
canonical comparison places two scenes beside each other: in Samuel the Lord
blesses Obed-edom's household; in Luke Elizabeth, filled with the Spirit,
pronounces Mary and the child blessed.  Luke does not say that Mary caused a
general blessing of Zechariah's house.

Then Luke gives the scene its long, quiet measure.  The Ark remained in the
house of Obed-edom for three months, and the Lord blessed the household
(2 Sam 6:11).  Mary remains with Elizabeth for about three months before
returning home (Luke 1:56).  The correspondence does not require the two women
to have occupied every day in a reconstructed schedule.  Its force lies in
the duration embedded within the same chain: journey through Judah, entrance
into a house, reverent question, embodied joy, the comparison of blessings,
and a three-month sojourn.

Luke states that Mary went with haste and remained about three months; he does
not narrate her daily activities.  Catholic contemplation may receive the
prompt journey and sustained presence as a form of charity: love is ready to
go and willing to remain, letting service acquire the ordinary weight of days.
That is a spiritual application, not a reconstructed diary of the Visitation.
It discloses a fitting order for apostolic life.  Presence is received before
it is carried, and what is carried is Christ, not the apostle's own importance.

\begin{studybox}{Stand at the threshold}
Read Exodus 40:34--35, 2 Samuel 6:9--11, and Luke 1:35, 39--56 without rushing
to an abstract conclusion.  Notice who acts, who receives, who is carried,
who rejoices, and who is blessed.  Then name the center of every scene: the
Lord chooses to dwell among his people.  Mary is the living Ark because the
Lord has truly made himself her Son.
\end{studybox}
